\documentclass[12pt]{article}

% PDF Layout
\usepackage[a4paper,left=2cm,right=5.5cm,bottom=3cm,top=2.5cm,marginparwidth=45mm,marginparsep=5mm]{geometry}

% Standard Packages
\usepackage[utf8]{inputenc}
\usepackage[T1]{fontenc}

\usepackage{hyperref}
\hypersetup{
	pdftitle={Note Template},
	pdfauthor={Euan J. Mendoza},
	colorlinks=true,
	linkcolor=black,
	citecolor=blue,
	pdfstartview=FitH,
	pdfpagemode=UseOutlines,
	pdfpagelayout=OneColumn,
}

% Personal Macros
\usepackage{em-mathtools}

\usepackage{todonotes}

% Bibliography Package
% \usepackage[style=alphabetic,maxbibnames=6]{biblatex}
% \addbibresource{references.bib}

\usepackage[capitalise]{cleveref}

% Useful for algorithms
% \usepackage{framed}

% Tikz setup

% \usepackage{standalone}
% \usepackage{tikz}

% \usetikzlibrary{cd}

\title{
    Note Template
}
\author{
    Euan Mendoza
}

\begin{document}

\maketitle

%% ============================================================
%% This file demonstrates features of the note template:
%%   - Wide Tufte-style right margin with \todo{} annotations
%%   - em-mathtools macros (number sets, groups, linear algebra, ...)
%%   - Theorem environments and cleveref cross-references
%% ============================================================

\section{Margin Annotations}

The wide right margin is for \todo{This is a todo note} inline reminders
and commentary using the \texttt{todonotes} package.
You can also place figures in the margin:\todo[inline]{Inline todos span the text width.}

\section{Number Sets and Fields}

The standard number sets are available as blackboard-bold commands:
\[
    \N \subset \Z \subset \Q \subset \R \subset \C.
\]
Generic fields and finite fields use $\F$ and $\K$.
For example, the $n$-dimensional vector space over $\F_{q}$ is $\F_{q}^{n}$.

\section{Group Theory}

\begin{definition}
    Let $X$ be a set. The \emph{symmetric group} on $X$ is
    \[
        \Sym(X) := \{ \sigma : X \to X \mid \sigma \text{ is a bijection} \}.
    \]
    The \emph{automorphism group}\todo{Generalises to any structured object.} of $X$ is the subgroup
    \[
        \Aut(X) := \{ \sigma \in \Sym(X) \mid \sigma \text{ preserves the structure of } X \}.
    \]
\end{definition}

\begin{definition}
    The \emph{general linear group} of degree $n$ over a field $\F$ is
    \[
        \GL(n, \F) := \{ A \in \M_{n}(\F) \mid \det A \neq 0 \}.
    \]
    The \emph{special orthogonal group} $\SO(n) \subset \GL(n, \R)$ consists of
    rotations in $\R^{n}$.
\end{definition}

\begin{theorem}[Cayley]
    \label{thm:cayley}
    Every group $G$ of order $n$ is isomorphic to a subgroup of $\Sym(\{1,\ldots,n\})$.
\end{theorem}

\begin{proof}
    Consider the left-regular action $G \to \Sym(G)$, $g \mapsto (x \mapsto gx)$.
    This is an injective group homomorphism, giving $G \cong \Iso(G)$ for some
    subgroup $\Iso(G) \leq \Sym(G)$.
\end{proof}

\section{Linear Algebra}

\begin{definition}
    The \emph{Grassmannian} $\Gr(k, \F^{n})$ is the set of all $k$-dimensional
    subspaces of $\F^{n}$.
\end{definition}

\begin{example}
    $\Gr(1, \F^{n}) = \mathbb{P}^{n-1}(\F)$ is projective space.
    A point $V \in \Gr(k, \F^{n})$ can be represented by a matrix
    $A \in \M_{k \times n}(\F)$ with $\rk(A) = k$, whose rows span $V$.
\end{example}

The \emph{inner product} on $\R^{n}$ is written $\ip{u}{v}$.
The \emph{transpose} of $A$ is $\tr{A}$, and its \emph{trace} is $\Tr(A)$.
The \emph{span} of vectors $v_{1}, \ldots, v_{k}$ is $\linspan\{v_{1}, \ldots, v_{k}\}$.

Bold letters denote vectors:\todo{Use $\vec{v}$ for column vectors, $\tr{\vec{v}}$ for row vectors.}
\[
    \vec{v} = \begin{pmatrix} v_{1} \\ \vdots \\ v_{n} \end{pmatrix}, \qquad
    \vzero = \begin{pmatrix} 0 \\ \vdots \\ 0 \end{pmatrix}.
\]

\section{Theorem Environments}

The following environments are provided by \texttt{em-mathtools}.
Cross-references use \texttt{cleveref}: see \cref{thm:cayley} for an example theorem,
and \cref{lem:example} below.

\begin{lemma}
    \label{lem:example}
    For any subspace $V \leq \F^{n}$,
    \[
        \dim V + \dim V^{\perp} = n.
    \]
\end{lemma}

\begin{proof}
    Apply the rank-nullity theorem to the map $\F^{n} \to V^{*}$, $x \mapsto \ip{\cdot}{x}|_{V}$.
\end{proof}

\begin{corollary}
    If $\dim V = k$ then $V^{\perp} \in \Gr(n-k, \F^{n})$.
\end{corollary}

\begin{remark}
    Over a general field $\F$, one must be careful: $V \cap V^{\perp}$ need not be trivial
    if $\F$ is not ordered.\todo{Consider $\F_{2}$: every vector is self-orthogonal.}
\end{remark}

\section{Probability and Complexity}

The expected value and variance are written $\E[X]$ and $\Var[X]$.
Common functions: $\poly(n)$, $\polylog(n)$.

The complexity class \NP\ contains \NPc\ problems.
\GI\ (Graph Isomorphism) is in \NP\ but not known to be \NPc.
It lies in \AM\ and \SZK.

\end{document}
